%% 1-13 九九の 表を たんけんしよう（10ページ、p.130〜139）
%% 第1巻のきまり：15 pt・行送り28 pt、分かち書き（文節の間に半角の空白）、漢字はすべてルビ、式は全角の × ＝ で書く
%% 1-12 と そろえる：「1つ分」（だん・左の かず・青）× 「いくつ分」（上の かず・だいだい）。おはじきは たて 1れつが 1つ分。

\makeatletter
% 囲みの 中の 文を、 右で そろえない（字の あいだが 広がらない ように）。 囲みの 中だけで きく
\newcommand\bkqlrag{\rightskip=0pt plus 1fil\relax}
% かんがえかたの 2行めを、 1行めの 書きだしに そろえる 字下げ（「かんがえかた」の 札の はば）
\newcommand\bkqlkindent{\setbox0=\hbox{\tikz\node[font=\sffamily\bfseries\small, inner xsep=5pt, inner ysep=2.5pt]{かんがえかた};}\hspace*{\dimexpr\wd0+0.6\zw\relax}}
% 九九の 表  \bkqlkuku{1マス(mm)}{色を つける 行(0=なし)}{色を つける 列}{0:見出しだけ 1:ぜんぶ 2:2〜5のだんだけ うすく 3:ななめだけ}{字の 大きさ}
%   左上が (0,0)。 マス (行 r, 列 c) の まん中は ((c+0.5)s, -(r+0.5)s)
\newcommand\bkqlkuku[5]{\begin{scope}
  \def\bkql@s{#1}%
  \ifnum#4=2 \def\bkql@tc{bkGray}\else\def\bkql@tc{bkInk}\fi
  \fill[white] (0,{-10*\bkql@s}) rectangle ({10*\bkql@s},0);
  \fill[bkTheme!12] (\bkql@s,0) rectangle ({10*\bkql@s},-\bkql@s);
  \fill[bkTheme!12] (0,-\bkql@s) rectangle (\bkql@s,{-10*\bkql@s});
  \ifnum#2>0 \fill[bkRoute2!18] (0,{-#2*\bkql@s}) rectangle ({(#3+1)*\bkql@s},{-(#2+1)*\bkql@s}); \fi
  \ifnum#3>0 \fill[bkOrange!20] ({#3*\bkql@s},0) rectangle ({(#3+1)*\bkql@s},{-(#2+1)*\bkql@s}); \fi
  \draw[bkGray!45, line width=0.45pt] (0,{-10*\bkql@s}) grid[step=\bkql@s] ({10*\bkql@s},0);
  \draw[bkInk!70, line width=0.9pt] (0,{-10*\bkql@s}) rectangle ({10*\bkql@s},0);
  \draw[bkInk!70, line width=0.9pt] (\bkql@s,0) -- (\bkql@s,{-10*\bkql@s}) (0,-\bkql@s) -- ({10*\bkql@s},-\bkql@s);
  \node[font=\sffamily\bfseries#5, text=bkTheme] at ({\bkql@s/2},{-\bkql@s/2}) {×};
  \foreach \bkql@i in {1,...,9} {
    \node[font=\sffamily\bfseries#5, text=bkOrange!85!black] at ({(\bkql@i+0.5)*\bkql@s},{-\bkql@s/2}) {\bkql@i};
    \node[font=\sffamily\bfseries#5, text=bkRoute2!85!black] at ({\bkql@s/2},{-(\bkql@i+0.5)*\bkql@s}) {\bkql@i};}
  \ifnum#4>0
  \foreach \bkql@r in {1,...,9} {\foreach \bkql@c in {1,...,9} {
    \pgfmathtruncatemacro\bkql@v{\bkql@r*\bkql@c}%
    \pgfmathtruncatemacro\bkql@show{(#4==1) || (#4==2 && \bkql@r>1 && \bkql@r<6) || (#4==3 && \bkql@r==\bkql@c)}%
    \ifnum\bkql@show=1 \node[font=#5, text=\bkql@tc, inner sep=0pt] at ({(\bkql@c+0.5)*\bkql@s},{-(\bkql@r+0.5)*\bkql@s}) {\bkql@v};\fi}}
  \fi
\end{scope}}
% おはじきの 四角（1-12 と 同じ）  \bkqlarray{(x,y)}{たての かず}{れつの かず}{間}{半径}{色}
\newcommand\bkqlarray[6]{\begin{scope}[shift={#1}]
  \foreach \bkql@i in {1,...,#3} {\foreach \bkql@j in {1,...,#2} {\bkohajiki{({(\bkql@i-1)*#4},{(\bkql@j-1)*#4})}{#5}{#6}}}
\end{scope}}
% たての 1れつ（1つ分）を かこむ わく  \bkqlcols{(x,y)}{たての かず}{れつの かず}{間}{半径}
\newcommand\bkqlcols[5]{\begin{scope}[shift={#1}]
  \foreach \bkql@i in {1,...,#3} {\draw[bkTheme!75, line width=0.7pt, dash pattern=on 2pt off 1.5pt, rounded corners=1.6mm]
    ({(\bkql@i-1)*#4-#5-1},{-#5-1}) rectangle ({(\bkql@i-1)*#4+#5+1},{(#2-1)*#4+#5+1});}
\end{scope}}
\makeatother

\bkchapter{1}{13}{\R{九九}{く|く}の \R{表}{ひょう}を\\たんけんしよう}{\bklevel{1}}{くくの ひょうを たんけんしよう}
\bkmichishirube{1-12}{1-14,1-17,1-21,7-1}

\begin{goal}
\begin{itemize}[label={\color{bkTheme}●}, leftmargin=1.4\zw, itemsep=0pt, topsep=0pt]
\item \R{九九}{く|く}の \R{表}{ひょう}で、 こたえを \R{見}{み}つけられる。
\item \R{表}{ひょう}の きまりを \R{見}{み}つけられる。
\item \R{入}{い}れかえても、 こたえは \R{同}{おな}じと わかる。
\end{itemize}
\end{goal}

\begin{hatena}\bkqlrag
\noindent\begin{minipage}[c]{0.56\linewidth}\centering
\begin{tikzpicture}[x=1mm, y=1mm]
  \fill[white, rounded corners=2mm] (-3,3) rectangle (68,-68);
  \bkqlkuku{6.5}{0}{0}{0}{\fontsize{9}{10}\selectfont}
  \foreach \k in {1,...,9} {\pgfmathtruncatemacro\v{\k*\k}
    \fill[bkYellow!55] ({\k*6.5+0.35},{-\k*6.5-0.35}) rectangle ({(\k+1)*6.5-0.35},{-(\k+1)*6.5+0.35});
    \node[font=\fontsize{9}{10}\selectfont\bfseries, inner sep=0pt] at ({(\k+0.5)*6.5},{-(\k+0.5)*6.5}) {\v};}
  \foreach \r/\c/\v/\col in {3/4/12/bkGreen,4/3/12/bkGreen,2/7/14/bkPurple,7/2/14/bkPurple} {
    \node[circle, fill=\col!20, draw=\col, line width=0.8pt, minimum size=5.6mm, inner sep=0pt, font=\fontsize{9}{10}\selectfont\bfseries] at ({(\c+0.5)*6.5},{-(\r+0.5)*6.5}) {\v};}
\end{tikzpicture}
\end{minipage}\hfill
\begin{minipage}[c]{0.42\linewidth}
ななめに 1、 4、 9、 16、 … と ならんで いるよ。 どんな かけ\R{算}{ざん}の こたえ かな？\par\vspace{4pt}
ななめの \R{線}{せん}を はさんで、 12と 12、 14と 14が むかいあって いるのは なぜ だろう？
\end{minipage}
\end{hatena}

\clearpage
%% ---------------- p.131 やってみよう ----------------
\begin{yattemiyou}[\R{九九}{く|く}の \R{表}{ひょう}を つくろう]\bkqlrag
\begin{enumerate}[label={\bkcircled[bkGreen]{\arabic*}}, leftmargin=2.2\zw, itemsep=0pt, topsep=0pt]
\item 2のだん〜5のだんは、 1-12で つくったね。
\item 1のだんを かこう。 1、 2、 3、 …
\item 6のだんは 6ずつ、 7のだんは 7ずつ ふえるよ（1-12）。 9のだんまで かこう。
\end{enumerate}
\begin{bkfigure}
\begin{tikzpicture}[x=1mm, y=1mm]
  \node[anchor=west, font=\sffamily\bfseries\footnotesize, text=bkRoute2!85!black] at (0,0) {6のだん};
  \foreach \i in {1,...,9} {\pgfmathtruncatemacro\v{6*\i}
    \node[circle, draw=bkGreen, line width=0.9pt, fill=white, minimum size=7.4mm, inner sep=0pt, font=\sffamily\bfseries\small, text=bkGreen!70!black] at ({14+\i*11.6},0) {\v};}
  \foreach \i in {1,...,8} {\draw[bk arrow, bkGreen, line width=0.8pt] ({14+\i*11.6+3.9},0) -- ({14+(\i+1)*11.6-3.9},0);
    \node[font=\sffamily\bfseries\fontsize{8}{9}\selectfont, text=bkGreen!70!black] at ({14+(\i+0.5)*11.6},4.6) {＋6};}
\end{tikzpicture}
\end{bkfigure}
\noindent{\small 6 ＋ 6 ＝ 12、 12 ＋ 6 ＝ 18、 … と、 まえの こたえに 6を たして いくよ。 10を こえる ときは、 1-7の くり\R{上}{あ}がりを つかおう。}
\begin{bkfigure}
\begin{tikzpicture}[x=1mm, y=1mm]
  \bkqlkuku{8.4}{0}{0}{2}{\small}
\end{tikzpicture}
\end{bkfigure}
\end{yattemiyou}

\bkfukidashi{mina}{ぜんぶ うめたら 81マス！ ゆっくりで だいじょうぶ。 まちがえても、 あとで なおせば いいよ。}

\clearpage
%% ---------------- p.132〜133 数学の道具 ----------------
\bksection{doguu}{\R{数学}{すう|がく}の \R{道具}{どう|ぐ}}

\bksubsection{1}{\R{表}{ひょう}の \R{見}{み}かた}

\label{bk:1-13:yomikata}\mbox{3 × 4} の こたえは、 \R{左}{ひだり}の 3から よこへ、 \R{上}{うえ}の 4から たてへ すすんで、 \R{出会}{で|あ}う マスに あります。 よこの ならびを 「\R{行}{ぎょう}」、 たての ならびを 「\R{列}{れつ}」と いいます。\bkidx{ぎょう}{\R{行}{ぎょう}（\R{表}{ひょう}の）}\bkidx{れつ}{\R{列}{れつ}（\R{表}{ひょう}の）}

\begin{bkfigure}
\begin{tikzpicture}[x=1mm, y=1mm]
  \bkqlkuku{8.4}{3}{4}{1}{\footnotesize}
  \draw[bkRed, line width=1.3pt] ({4.5*8.4},{-3.5*8.4}) circle (3.9);
  \draw[bk arrow, bkRoute2, line width=1.1pt] (88.5,{-3.5*8.4}) -- (85,{-3.5*8.4});
  \node[anchor=west, font=\sffamily\bfseries\footnotesize, text=bkRoute2!85!black, align=left] at (89,{-3.5*8.4}) {3の \R{行}{ぎょう}（よこ）\\1つ\R{分}{ぶん}の かず};
  \draw[bk arrow, bkOrange, line width=1.1pt] ({4.5*8.4},6.2) -- ({4.5*8.4},0.6);
  \node[anchor=south, font=\sffamily\bfseries\footnotesize, text=bkOrange!85!black] at ({4.5*8.4},6) {4の \R{列}{れつ}（たて）\ いくつ\R{分}{ぶん}};
  \node[anchor=west, font=\sffamily\bfseries\small, text=bkRed] at (89,-60) {3 × 4 ＝ 12};
\end{tikzpicture}
\end{bkfigure}
\noindent どうして \R{出会}{で|あ}う マスなのかな？ 3の \R{行}{ぎょう}には、 3のだん（\mbox{3 × 1}、 \mbox{3 × 2}、 …）が ならんで います。 4の \R{列}{れつ}には、 「× 4」の こたえが ならんで います。 その どちらにも \R{入}{はい}って いる マスが、 \mbox{3 × 4} なのです。

\bksubsection{2}{ふえかたの きまり}

\label{bk:1-13:fuekata}3のだんは 3ずつ ふえます。 いくつ\R{分}{ぶん}が 1 ふえると、 こたえは 1つ\R{分}{ぶん}の かずだけ ふえます。

\begin{bkfigure}
\begin{tikzpicture}[x=1mm, y=1mm]
  \fill[bkRoute2!18] (0,0) rectangle (18,8);
  \node[font=\sffamily\bfseries\footnotesize, text=bkRoute2!85!black] at (9,4) {3のだん};
  \foreach \i in {1,...,9} {\pgfmathtruncatemacro\v{3*\i}
    \node[font=\small] at ({18+(\i-0.5)*11},4) {\v};}
  \foreach \i in {1,...,8} {\draw[bkGray!45, line width=0.45pt] ({18+\i*11},0) -- ({18+\i*11},8);}
  \draw[bkInk!70, line width=0.9pt] (0,0) rectangle (117,8) (18,0) -- (18,8);
  \foreach \i in {1,...,8} {\draw[bkTheme, line width=1pt, -{Stealth[length=4pt]}] ({18+(\i-0.5)*11+1.5},8.8) .. controls ({18+(\i-0.5)*11+4},12.5) and ({18+(\i+0.5)*11-4},12.5) .. ({18+(\i+0.5)*11-1.5},8.8);
    \node[font=\sffamily\bfseries\fontsize{8}{9}\selectfont, text=bkTheme] at ({18+\i*11},14.3) {＋3};}
\end{tikzpicture}
\end{bkfigure}
\noindent こんどは、 \R{表}{ひょう}を たてに \R{見}{み}よう。 4の \R{列}{れつ}は \mbox{1 × 4 ＝ 4}、 \mbox{2 × 4 ＝ 8}、 \mbox{3 × 4 ＝ 12}、 … と、 4ずつ ふえます。 なぜ 4 なのかな？
\begin{bkfigure}
\begin{tikzpicture}[x=1mm, y=1mm]
  \bkqlcols{(6,0)}{3}{4}{7.5}{2.7}
  \bkqlarray{(6,0)}{3}{4}{7.5}{2.7}{bkRoute2}
  \node[font=\sffamily\bfseries\footnotesize] at (17.25,-7) {3 × 4 ＝ 12};
  \draw[bk arrow, bkTheme, line width=1.2pt] (40,8) -- (52,8);
  \bkqlcols{(64,0)}{4}{4}{7.5}{2.7}
  \bkqlarray{(64,0)}{3}{4}{7.5}{2.7}{bkRoute2}
  \foreach \i in {0,...,3} {\bkohajiki{({64+\i*7.5},22.5)}{2.7}{bkOrange}}
  \node[font=\sffamily\bfseries\footnotesize] at (75.25,-7) {4 × 4 ＝ 16};
  \node[anchor=west, font=\sffamily\bfseries\footnotesize, text=bkOrange!85!black, align=left] at (93,22.5) {どの れつにも\\1こずつ ふえた};
  \node[anchor=west, font=\sffamily\bfseries\footnotesize, text=bkOrange!85!black] at (93,10) {4れつで ＋4};
\end{tikzpicture}
\end{bkfigure}
\noindent 1つ\R{分}{ぶん}が 3から 4に ふえると、 4れつ どれにも 1こずつ ふえます。 4れつ あるので、 ぜんぶで 4 ふえるのです。

\noindent\begin{minipage}[c]{0.6\linewidth}
9のだんにも おもしろい きまりが あります。 \R{十}{じゅう}の くらいは 1ずつ ふえ、 \R{一}{いち}の くらいは 1ずつ へります。 9を たすのは、 「10 たして 1 ひく」のと \R{同}{おな}じ だからです。 \mbox{18 ＋ 9} なら、 \mbox{18 ＋ 10 ＝ 28}、 \mbox{28 － 1 ＝ 27}。
\end{minipage}\hfill
\begin{minipage}[c]{0.36\linewidth}\centering
\begin{tblr}{colspec={Q[c,m,wd=12mm]Q[c,m,wd=9mm]Q[c,m,wd=9mm]}, hlines={0.5pt,bkLine}, vlines={0.5pt,bkLine}, rowsep=0pt, colsep=1pt, cells={font=\footnotesize}, row{1}={bg=bkTheme!12, font=\sffamily\bfseries\footnotesize}}
 & \R{十}{じゅう} & \R{一}{いち} \\
9 × 1 & 0 & 9 \\
9 × 2 & 1 & 8 \\
9 × 3 & 2 & 7 \\
9 × 4 & 3 & 6 \\
9 × 5 & 4 & 5 \\
\end{tblr}
\end{minipage}

\bksubsection{3}{\R{入}{い}れかえても \R{同}{おな}じ}

\label{bk:1-13:irekae}\mbox{3 × 4} の おはじきの \mbox{\R{四角}{し|かく}}を、 くるっと まわすと \mbox{4 × 3} の \mbox{\R{四角}{し|かく}}に なります。 おはじきの かずは かわりません。

\begin{bkfigure}
\begin{tikzpicture}[x=1mm, y=1mm]
  \bkqlcols{(10,4)}{3}{4}{7.5}{2.7}
  \bkqlarray{(10,4)}{3}{4}{7.5}{2.7}{bkRoute2}
  \node[font=\sffamily\bfseries\small] at (21.25,-6) {3 × 4 ＝ 12};
  \draw[bkTheme, line width=1.4pt, -{Stealth[length=7pt]}] (46,24) arc (120:60:16) node[midway, above=0.5mm, font=\sffamily\bfseries\footnotesize, text=bkTheme] {くるっと};
  \bkqlcols{(84,0)}{4}{3}{7.5}{2.7}
  \bkqlarray{(84,0)}{4}{3}{7.5}{2.7}{bkRoute4}
  \node[font=\sffamily\bfseries\small] at (91.5,-6) {4 × 3 ＝ 12};
\end{tikzpicture}
\end{bkfigure}

\begin{taisetsu}
かけ\R{算}{ざん}は、 \R{入}{い}れかえても こたえは \R{同}{おな}じ。\\
\mbox{3 × 4 ＝ 4 × 3}
\end{taisetsu}\bkidx{いれかえ}{\R{入}{い}れかえても \R{同}{おな}じ}

\noindent\begin{minipage}[c]{0.55\linewidth}
この きまりを つかうと、 おぼえる ことが へります。 \mbox{7 × 3} を わすれても、 \mbox{3 × 7 ＝ 21} を おもいだせば いいのです。 \R{右}{みぎ}の \R{表}{ひょう}の \R{色}{いろ}の ついた マス（45マス）だけ おぼえれば、 のこりは \R{入}{い}れかえて わかります。
\end{minipage}\hfill
\begin{minipage}[c]{0.42\linewidth}\centering
\begin{tikzpicture}[x=1mm, y=1mm]
  \foreach \r in {1,...,9} {\foreach \c in {1,...,9} {
    \ifnum\c<\r \def\bkqlcol{white}\else\def\bkqlcol{bkYellow!45}\fi
    \filldraw[fill=\bkqlcol, draw=bkGray!45, line width=0.4pt] ({(\c-1)*5.2},{-(\r-1)*5.2}) rectangle ++(5.2,-5.2);}}
  \draw[bkInk!70, line width=0.8pt] (0,0) rectangle (46.8,-46.8);
  \foreach \k in {1,...,9} {\node[font=\sffamily\bfseries\fontsize{7}{8}\selectfont, text=bkOrange!85!black] at ({(\k-0.5)*5.2},2.6) {\k};
    \node[font=\sffamily\bfseries\fontsize{7}{8}\selectfont, text=bkRoute2!85!black] at (-2.6,{-(\k-0.5)*5.2}) {\k};}
\end{tikzpicture}
\end{minipage}

\bksubsection{4}{ななめの きまり}

\label{bk:1-13:naname}ななめの かずは、 \mbox{1 × 1}、 \mbox{2 × 2}、 \mbox{3 × 3}、 … \R{同}{おな}じ かずどうしの かけ\R{算}{ざん}。 おはじきは ましかくに ならびます。

\begin{bkfigure}
\begin{tikzpicture}[x=1mm, y=1mm]
  \foreach \n/\c in {1/12,2/42,3/75,4/110} {
    \bkqlarray{({\c-(\n-1)*2.75},0)}{\n}{\n}{5.5}{2.1}{bkRoute5}
    \pgfmathtruncatemacro\a{\n*\n}
    \node[font=\sffamily\bfseries\footnotesize, text=bkInk] at (\c,-6.5) {\n × \n ＝ \a};}
\end{tikzpicture}
\end{bkfigure}
\noindent \bkhatenaref の 12と 12は、 \mbox{3 × 4} と \mbox{4 × 3}。 ③の きまりで \R{同}{おな}じ かずだから、 ななめの \R{線}{せん}を はさんで むかいあうよ。

\bksubsection{5}{だんを たして つくる}

\label{bk:1-13:tasu}\R{表}{ひょう}の 2のだんと 3のだんを、 たてに たして みよう。
\begin{bkfigure}
\begin{tblr}{colspec={Q[c,m,wd=22mm]*{6}{Q[c,m,wd=11mm]}Q[c,m,wd=8mm]}, hlines={0.6pt,bkLine}, vlines={0.6pt,bkLine},
  column{1}={font=\sffamily\bfseries\small}, rowsep=1pt, colsep=0pt, cells={font=\small}, row{3}={bg=bkRoute2!14}, row{4}={bg=bkYellow!40}}
 & × 1 & × 2 & × 3 & × 4 & × 5 & × 6 & … \\
2のだん & 2 & 4 & 6 & 8 & 10 & 12 & … \\
3のだん & 3 & 6 & 9 & 12 & 15 & 18 & … \\
たすと & 5 & 10 & 15 & 20 & 25 & 30 & … \\
\end{tblr}
\end{bkfigure}
\noindent\begin{minipage}[c]{0.6\linewidth}
たすと、 5のだんに なりました。 おはじきで \R{見}{み}ると、 1れつが 2こと 3こで 5こ。 だから \mbox{2 × 4} と \mbox{3 × 4} を たすと \mbox{5 × 4} です。\par
むずかしい \R{九九}{く|く}も、 これで もとめられます。 \mbox{6 × 7} なら、 6を 5と 1に わけて \mbox{5 × 7 ＝ 35}、 \mbox{1 × 7 ＝ 7}。 \mbox{35 ＋ 7 ＝ 42} です。
\end{minipage}\hfill
\begin{minipage}[c]{0.36\linewidth}\centering
\begin{tikzpicture}[x=1mm, y=1mm]
  \foreach \i in {0,...,3} {\foreach \j in {0,1} {\bkohajiki{(\i*7.5,\j*6.5)}{2.6}{bkRoute2}}
    \foreach \j in {2,3,4} {\bkohajiki{(\i*7.5,\j*6.5+2.5)}{2.6}{bkOrange}}
    \draw[bkTheme!75, line width=0.7pt, dash pattern=on 2pt off 1.5pt, rounded corners=1.6mm] ({\i*7.5-3.8},-3.8) rectangle ({\i*7.5+3.8},32.3);}
  \node[anchor=west, font=\sffamily\bfseries\footnotesize, text=bkRoute2!85!black] at (27,3.25) {2 × 4};
  \node[anchor=west, font=\sffamily\bfseries\footnotesize, text=bkOrange!85!black] at (27,21.5) {3 × 4};
  \node[font=\sffamily\bfseries\footnotesize] at (11.25,-8) {5 × 4 ＝ 20};
\end{tikzpicture}
\end{minipage}


\begin{yattemiyou}[7のだんを 5のだんと 2のだんで つくろう]\bkqlrag
5のだんと 2のだんを たてに たして、 7のだんを かこう。
\par\vspace{3pt}
\noindent\hfill\begin{tblr}{colspec={Q[c,m,wd=18mm]*{9}{Q[c,m,wd=10.5mm]}}, hlines={0.6pt,bkLine}, vlines={0.6pt,bkLine},
  column{1}={font=\sffamily\bfseries\small}, rowsep=2pt, colsep=0pt, cells={font=\small}, row{1}={bg=bkTheme!12, font=\sffamily\bfseries\small}}
 & × 1 & × 2 & × 3 & × 4 & × 5 & × 6 & × 7 & × 8 & × 9 \\
5のだん & 5 & 10 & 15 & 20 & 25 & 30 & 35 & 40 & 45 \\
2のだん & 2 & 4 & 6 & 8 & 10 & 12 & 14 & 16 & 18 \\
7のだん & 7 & \bkkakikomi[8.5mm] & \bkkakikomi[8.5mm] & \bkkakikomi[8.5mm] & \bkkakikomi[8.5mm] & \bkkakikomi[8.5mm] & \bkkakikomi[8.5mm] & \bkkakikomi[8.5mm] & \bkkakikomi[8.5mm] \\
\end{tblr}\hfill\null
\par\vspace{3pt}
できたら、 やってみようで つくった \R{九九}{く|く}の \R{表}{ひょう}の 7のだんと くらべよう。 \R{同}{おな}じに なったかな？
\end{yattemiyou}

\clearpage
%% ---------------- p.134〜135 例題 ----------------

\begin{reidai}[\R{表}{ひょう}で \R{見}{み}つける]
\R{九九}{く|く}の \R{表}{ひょう}で、 \mbox{7 × 6} の こたえを \R{見}{み}つけよう。

\noindent\begin{minipage}[c]{0.7\linewidth}
\kangaekata 7の \R{行}{ぎょう}と 6の \R{列}{れつ}。

\shiki 7 × 6 ＝ 42 \kotaenext 42
\end{minipage}\hfill
\begin{minipage}[c]{0.28\linewidth}\centering
\begin{tikzpicture}[x=1mm, y=1mm]
  \begin{scope}[scale=0.85]
  \foreach \c/\t in {1/5,2/6,3/7} {\fill[bkTheme!12] ({\c*7},0) rectangle ({(\c+1)*7},-7);
    \node[font=\sffamily\bfseries\footnotesize, text=bkOrange!85!black] at ({(\c+0.5)*7},-3.5) {\t};}
  \foreach \r/\t in {1/6,2/7,3/8} {\fill[bkTheme!12] (0,{-\r*7}) rectangle (7,{-(\r+1)*7});
    \node[font=\sffamily\bfseries\footnotesize, text=bkRoute2!85!black] at (3.5,{-(\r+0.5)*7}) {\t};}
  \fill[bkRoute2!18] (7,-14) rectangle (21,-21);
  \fill[bkOrange!20] (14,-7) rectangle (21,-14);
  \foreach \r/\a in {1/6,2/7,3/8} {\foreach \c/\b in {1/5,2/6,3/7} {\pgfmathtruncatemacro\v{\a*\b}
    \node[font=\footnotesize] at ({(\c+0.5)*7},{-(\r+0.5)*7}) {\v};}}
  \draw[bkGray!45, line width=0.45pt] (0,-28) grid[step=7] (28,0);
  \draw[bkInk!70, line width=0.8pt] (0,-28) rectangle (28,0) (7,0) -- (7,-28) (0,-7) -- (28,-7);
  \draw[bkRed, line width=1.2pt] (17.5,-17.5) circle (3.3);
  \end{scope}
\end{tikzpicture}
\end{minipage}
\end{reidai}

\begin{jibun}
\R{九九}{く|く}の \R{表}{ひょう}で、 こたえを \R{見}{み}つけよう。\par\vspace{2pt}
\noindent\hspace*{1\zw}（1）\mbox{8 × 4 ＝} \bkkakikomi\hspace{4\zw}（2）\mbox{9 × 7 ＝} \bkkakikomi
\bkjans{（1）32\quad（2）63}
\end{jibun}

\begin{reidai}[ふえかたを つかう]
あめが 6こずつ 7ふくろで 42こ。 8ふくろなら なんこ？

\kangaekata 1ふくろ ふえると 6こ ふえる。\\ \mbox{42 ＋ 6 ＝ 48}

\shiki 6 × 8 ＝ 48 \kotaenext 48こ
\end{reidai}

\begin{jibun}
クレヨンが 7\R{本}{ほん}ずつ 5はこで 35\R{本}{ぼん}。 6はこなら なん\R{本}{ぼん}？
\bkshikiran
\bkjans{\mbox{7 × 6 ＝ 42}\quad 42\R{本}{ほん}}
\end{jibun}

\begin{reidai}[\R{入}{い}れかえて けいさん]
\mbox{7 × 3} の こたえを、 3のだんで もとめよう。

\kangaekata \mbox{7 × 3} は \mbox{3 × 7} と \R{同}{おな}じ。 3のだんで 21。

\shiki 7 × 3 ＝ 21 \kotaenext 21
\end{reidai}

\begin{jibun}
（1）\mbox{9 × 4} の こたえを、 4のだんで もとめよう。\hfill\bkkakikomi\par\vspace{2pt}
（2）□に あう かずを かこう。\hspace{2\zw}\mbox{6 × 5 ＝ 5 ×} \bkkakikomi
\bkjans{（1）36\quad（2）6}
\end{jibun}

\begin{reidai}[こたえから さがす]
\mbox{3 × □ ＝ 12} の □に \R{入}{はい}る かずを、 \R{表}{ひょう}で さがそう。

\kangaekata 3の \R{行}{ぎょう}を \R{右}{みぎ}へ。 3、 6、 9、 12。\\ 12の マスは、 4の \R{列}{れつ}に あるね。

\kotae 4\par
{\small\color{bkBlue!80!black}この さがしかたは、 1-14の わり\R{算}{ざん}で つかうよ。}
\end{reidai}

\begin{jibun}
（1）□に \R{入}{はい}る かずを、 \R{表}{ひょう}で さがそう。\hfill \mbox{4 ×} \bkkakikomi\enspace ＝ 28\par\vspace{2pt}
（2）□に \R{入}{はい}る かずを、 \R{表}{ひょう}で さがそう。\hfill \mbox{6 ×} \bkkakikomi\enspace ＝ 54
\bkjans{（1）7\quad（2）9}
\end{jibun}

\begin{reidai}[だんを たして もとめる]
\mbox{7 × 8} を、 5のだんと 2のだんで もとめよう。

\kangaekata 7は 5と 2。 \mbox{5 × 8 ＝ 40}、 \mbox{2 × 8 ＝ 16}。\\
\bkqlkindent\mbox{40 ＋ 16 ＝ 56}

\shiki 7 × 8 ＝ 56 \kotaenext 56
\end{reidai}

\begin{jibun}
\mbox{8 × 6} を、 5のだんと 3のだんで もとめよう。\par\vspace{2pt}
\noindent\hspace*{1\zw}\mbox{5 × 6 ＝} \bkkakikomi\enspace\mbox{3 × 6 ＝} \bkkakikomi\enspace\mbox{8 × 6 ＝} \bkkakikomi
\bkjans{\mbox{5 × 6 ＝ 30}、 \mbox{3 × 6 ＝ 18}、 \mbox{30 ＋ 18 ＝ 48}\quad \mbox{8 × 6 ＝ 48}}
\end{jibun}

\begin{reidai}[こたえが \R{同}{おな}じ マスを ぜんぶ さがす]
\R{九九}{く|く}の \R{表}{ひょう}で、 こたえが 18の マスを ぜんぶ さがそう。

\kangaekata 1のだんから じゅんに、 1だんずつ しらべる。\\
\bkqlkindent 1のだん：ない（9までしか ない）\\
\bkqlkindent 2のだん：\mbox{2 × 9}\quad 3のだん：\mbox{3 × 6}\\
\bkqlkindent 4・5のだん：ない\quad 6のだん：\mbox{6 × 3}\\
\bkqlkindent 7・8のだん：ない\quad 9のだん：\mbox{9 × 2}

\kotae \mbox{2 × 9}、 \mbox{3 × 6}、 \mbox{6 × 3}、 \mbox{9 × 2} の 4つ
\end{reidai}

\begin{jibun}
こたえが 12の マスは、 いくつ ある？ 1のだんから じゅんに しらべよう。\hfill\bkkakikomi\enspace こ
\bkjans{4こ（\mbox{2 × 6}、 \mbox{3 × 4}、 \mbox{4 × 3}、 \mbox{6 × 2}）}
\end{jibun}

\begin{reidai}[\R{表}{ひょう}の そとへ]
\mbox{3 × 10} は いくつ？ \R{表}{ひょう}には ないけれど、 ふえかたの きまりで もとめよう。

\kangaekata \mbox{3 × 9 ＝ 27}。 いくつ\R{分}{ぶん}が 1 ふえると、 3 ふえる。\\
\bkqlkindent\mbox{27 ＋ 3 ＝ 30}

\shiki 3 × 10 ＝ 30 \kotaenext 30
\end{reidai}

\begin{jibun}
\mbox{4 × 10} は いくつ？ \mbox{4 × 9 ＝ 36} を つかおう。
\bkshikiran[3\zw]
\bkjans{\mbox{36 ＋ 4 ＝ 40}\quad \mbox{4 × 10 ＝ 40}}
\end{jibun}

\begin{reidai}[\R{表}{ひょう}の いちぶを うめる]
\noindent\begin{minipage}[c]{0.6\linewidth}
\R{九九}{く|く}の \R{表}{ひょう}の いちぶです。 ア、 イに \R{入}{はい}る かずは？

\kangaekata \R{左}{ひだり} × \R{上}{うえ}\\
\bkqlkindent ア は \mbox{4 × 6 ＝ 24}\\
\bkqlkindent イ は \mbox{5 × 7 ＝ 35}

\kotae ア 24\quad イ 35
\end{minipage}\hfill
\begin{minipage}[c]{0.38\linewidth}\centering
\begin{tikzpicture}[x=1mm, y=1mm]
  \def\s{9}
  \fill[bkTheme!12] (\s,0) rectangle ({3*\s},-\s) (0,-\s) rectangle (\s,{-3*\s});
  \foreach \c/\t in {1/6,2/7} {\node[font=\sffamily\bfseries\small, text=bkOrange!85!black] at ({(\c+0.5)*\s},{-\s/2}) {\t};}
  \foreach \r/\t in {1/4,2/5} {\node[font=\sffamily\bfseries\small, text=bkRoute2!85!black] at ({\s/2},{-(\r+0.5)*\s}) {\t};}
  \node[font=\small] at ({2.5*\s},{-1.5*\s}) {28};
  \node[font=\small] at ({1.5*\s},{-2.5*\s}) {30};
  \node[font=\sffamily\bfseries\small, text=bkBlue] at ({1.5*\s},{-1.5*\s}) {ア};
  \node[font=\sffamily\bfseries\small, text=bkBlue] at ({2.5*\s},{-2.5*\s}) {イ};
  \draw[bkGray!45, line width=0.45pt] (0,{-3*\s}) grid[step=\s] ({3*\s},0);
  \draw[bkInk!70, line width=0.9pt] (0,{-3*\s}) rectangle ({3*\s},0) (\s,0) -- (\s,{-3*\s}) (0,-\s) -- ({3*\s},-\s);
  \node[font=\sffamily\bfseries\small, text=bkTheme] at ({\s/2},{-\s/2}) {×};
\end{tikzpicture}
\end{minipage}
\end{reidai}

\begin{jibun}
\noindent\begin{minipage}[c]{0.6\linewidth}
ウ、 エに \R{入}{はい}る かずを かこう。\par\vspace{4pt}
\hspace*{1\zw}ウ\enspace\bkkakikomi\hspace{2\zw}エ\enspace\bkkakikomi
\bkjans{ウ \mbox{3 × 9 ＝ 27}\quad エ \mbox{4 × 8 ＝ 32}}
\end{minipage}\hfill
\begin{minipage}[c]{0.38\linewidth}\centering
\begin{tikzpicture}[x=1mm, y=1mm]
  \def\s{9}
  \fill[bkTheme!12] (\s,0) rectangle ({3*\s},-\s) (0,-\s) rectangle (\s,{-3*\s});
  \foreach \c/\t in {1/8,2/9} {\node[font=\sffamily\bfseries\small, text=bkOrange!85!black] at ({(\c+0.5)*\s},{-\s/2}) {\t};}
  \foreach \r/\t in {1/3,2/4} {\node[font=\sffamily\bfseries\small, text=bkRoute2!85!black] at ({\s/2},{-(\r+0.5)*\s}) {\t};}
  \node[font=\small] at ({1.5*\s},{-1.5*\s}) {24};
  \node[font=\small] at ({2.5*\s},{-2.5*\s}) {36};
  \node[font=\sffamily\bfseries\small, text=bkOrange] at ({2.5*\s},{-1.5*\s}) {ウ};
  \node[font=\sffamily\bfseries\small, text=bkOrange] at ({1.5*\s},{-2.5*\s}) {エ};
  \draw[bkGray!45, line width=0.45pt] (0,{-3*\s}) grid[step=\s] ({3*\s},0);
  \draw[bkInk!70, line width=0.9pt] (0,{-3*\s}) rectangle ({3*\s},0) (\s,0) -- (\s,{-3*\s}) (0,-\s) -- ({3*\s},-\s);
  \node[font=\sffamily\bfseries\small, text=bkTheme] at ({\s/2},{-\s/2}) {×};
\end{tikzpicture}
\end{minipage}
\end{jibun}

%% ---------------- p.136 まちがえやすいところ ----------------
\bksection{tsumazuki}{まちがえやすい ところ}

\begin{tsumazuki}{となりの マスを よむ}\bkqlrag
\bkbatsu\enspace \mbox{7 × 6 ＝ 48}\hspace{3\zw}\bkmaru\enspace \mbox{7 × 6 ＝ 42}\par
ゆびを 2\R{本}{ほん} つかって、 \R{行}{ぎょう}と \R{列}{れつ}を たどろう。
\end{tsumazuki}

\begin{tsumazuki}{にて いる こたえ}\bkqlrag
\bkbatsu\enspace \mbox{7 × 8 ＝ 54}\hspace{3\zw}\bkmaru\enspace \mbox{7 × 8 ＝ 56}\par
まよったら、 \mbox{7 × 7 ＝ 49} に 7を たして たしかめよう。 49は あと 1で 50。 のこりの 6で 56（1-4の 「10の まとまり」と おなじ）。
\end{tsumazuki}

\begin{tsumazuki}{ふえかたを まちがえる}\bkqlrag
\mbox{6 × 7 ＝ 42}。 では \mbox{6 × 8} は？\par
\bkbatsu\enspace \mbox{42 ＋ 1 ＝ 43}\hspace{3\zw}\bkmaru\enspace \mbox{42 ＋ 6 ＝ 48}\par
いくつ\R{分}{ぶん}が 1 ふえると、 1つ\R{分}{ぶん}の 6が ふえるよ。
\end{tsumazuki}

\begin{column}{ななめの かずの ふしぎ}\bkqlrag
\noindent\begin{minipage}[c]{0.62\linewidth}
1、 4、 9、 16 の ふえかたを \R{見}{み}よう。 3、 5、 7 と、 2ずつ ふえて いるね。 ましかくの まわりに、 L の かたちを かさねて いくと わかるよ。
\end{minipage}\hfill
\begin{minipage}[c]{0.36\linewidth}\centering
\begin{tikzpicture}[x=1mm, y=1mm]
  \foreach \i in {1,...,4} {\foreach \j in {1,...,4} {
    \pgfmathtruncatemacro\k{max(\i,\j)}
    \ifcase\k\or\def\col{bkRoute5}\or\def\col{bkRoute2}\or\def\col{bkRoute4}\or\def\col{bkOrange}\fi
    \bkohajiki{({(\i-1)*6},{(4-\j)*6})}{2.3}{\col}}}
  \foreach \k/\t/\col in {1/1/bkRoute5,2/＋3/bkRoute2,3/＋5/bkRoute4,4/＋7/bkOrange} {
    \node[anchor=east, font=\sffamily\bfseries\footnotesize, text=\col!85!black] at (-3.5,{(4-\k)*6+0.4}) {\t};}
\end{tikzpicture}
\end{minipage}
\end{column}

\clearpage
%% ---------------- p.137 つかわれているところ ----------------
\begin{butsuri}\bkqlrag
\begin{bkfigure}
\begin{tikzpicture}[x=1mm, y=1mm]
  % 時間わり
  \fill[white, rounded corners=3mm] (0,0) rectangle (73,42);
  \begin{scope}[shift={(2.5,37)}]
    \foreach \c/\t in {1/げつ,2/か,3/すい,4/もく,5/きん} {\fill[bkTheme!12] ({7+(\c-1)*12.2},0) rectangle ({7+\c*12.2},-6);
      \node[font=\sffamily\bfseries\fontsize{8}{9}\selectfont] at ({7+(\c-0.5)*12.2},-3) {\t};}
    \foreach \r in {1,2,3} {\fill[bkTheme!12] (0,{-6-(\r-1)*7}) rectangle (7,{-6-\r*7});
      \node[font=\sffamily\bfseries\fontsize{8}{9}\selectfont] at (3.5,{-6-(\r-0.5)*7}) {\r};}
    \fill[bkYellow!45] (31.4,-13) rectangle (43.6,-20);
    \foreach \r/\row in {1/{こくご,さんすう,こくご,さんすう,たいいく},2/{さんすう,こくご,さんすう,おんがく,こくご},3/{せいかつ,ずこう,たいいく,こくご,せいかつ}} {
      \foreach \t [count=\c] in \row {\node[font=\fontsize{7}{8}\selectfont, inner sep=0pt] at ({7+(\c-0.5)*12.2},{-6-(\r-0.5)*7}) {\t};}}
    \foreach \c in {1,...,5} \draw[bkGray!40, line width=0.4pt] ({7+(\c-1)*12.2},0) -- ({7+(\c-1)*12.2},-27);
    \foreach \r in {1,2,3} \draw[bkGray!40, line width=0.4pt] (0,{-6-(\r-1)*7}) -- (68,{-6-(\r-1)*7});
    \draw[bkInk!60, line width=0.8pt] (0,0) rectangle (68,-27);
  \end{scope}
  \node[font=\sffamily\bfseries\small] at (36.5,4.5) {\R{時間}{じ|かん}わり};
  % がめんの 点
  \begin{scope}[shift={(77,0)}]
  \fill[white, rounded corners=3mm] (0,0) rectangle (45,42);
  \fill[bkInk!85, rounded corners=1.5mm] (3,8.5) rectangle (42,39.5);
  \begin{scope}[shift={(6.5,12)}]
    \foreach \row [count=\r] in {{0,1,1,0,0,1,1,0},{1,1,1,1,1,1,1,1},{1,1,1,1,1,1,1,1},{1,1,1,1,1,1,1,1},{0,1,1,1,1,1,1,0},{0,0,1,1,1,1,0,0},{0,0,0,1,1,0,0,0}} {
      \foreach \v [count=\c] in \row {\ifnum\v=1 \fill[bkAppleRed] ({(\c-1)*4},{(7-\r)*4}) rectangle ++(3.6,3.6);
        \else \fill[bkInk!65] ({(\c-1)*4},{(7-\r)*4}) rectangle ++(3.6,3.6);\fi}}
  \end{scope}
  \node[font=\sffamily\bfseries\small] at (22.5,4.5) {\R{画面}{が|めん}の \R{点}{てん}};
  \end{scope}
\end{tikzpicture}
\end{bkfigure}
\begin{itemize}[label={\color{bkPurple}●}, leftmargin=1.4\zw, itemsep=0pt, topsep=0pt]
\item \R{時間}{じ|かん}わりや \R{電車}{でん|しゃ}の \R{時}{じ}こく\R{表}{ひょう}も \R{表}{ひょう}。 「すいようびの 2\R{時間}{じ|かん}\R{目}{め}」は、 \R{行}{ぎょう}と \R{列}{れつ}が \R{出会}{で|あ}う ところ。
\item テレビや ゲームの \R{画面}{が|めん}は、 \R{小}{ちい}さな \R{点}{てん}が たて・よこに ならんだ \R{表}{ひょう}です。
\item ロボットの うでの うごきも、 コンピュータが かずの \R{表}{ひょう}で けいさんして います。
\end{itemize}
\end{butsuri}

\bkfukidashi{delta}{たて・よこの \R{表}{ひょう}で けいさんを まとめる ことを 「\R{行列}{ぎょう|れつ}」と いうよ。 \R{第}{だい}7\R{巻}{かん}で まなぶ、 ロボットや コンピュータの だいじな \R{道具}{どう|ぐ}なんだ！}

\begin{bkfigure}
\begin{tikzpicture}[x=1mm, y=1mm]
  \draw[bkRoute5, line width=3pt] (10,0) -- (128,0);
  \node[anchor=west, inner sep=0pt] at (0,12) {\bkrouteicon[1.4em]{5}\hspace{0.3em}{\sffamily\bfseries\small\color{bkRoute5!80!black}\R{表}{ひょう}の \R{路線}{ろ|せん}}};
  \foreach \x/\c/\t/\s in {10/bkVol1/{1-13}/{\R{九九}{く|く}の \R{表}{ひょう}},49.3/bkVol1/{1-17}/{\R{表}{ひょう}と グラフ},88.7/bkVol3/{3-6}/{れんりつ},128/bkVol7/{7-1}/{\R{行列}{ぎょう|れつ}}} {
    \fill[white, draw=\c, line width=1.6pt] (\x,0) circle (2.6);
    \node[anchor=base, font=\sffamily\bfseries\footnotesize, text=\c!85!black] at (\x,-8.4) {\t};
    \node[anchor=base, font=\sffamily\footnotesize, text=bkInk] at (\x,-15.4) {\s};}
  \node[fill=bkTheme, text=white, rounded corners=1mm, font=\sffamily\bfseries\fontsize{8}{9}\selectfont, inner sep=1.5pt] at (10,5.5) {いま ここ};
\end{tikzpicture}
\end{bkfigure}

\clearpage
%% ---------------- p.138〜139 もんだい ----------------
\bkendproblems

\bkproblemgroup{きほん}{1}

\bkproblem けいさんを しましょう。
\par\vspace{2pt}
\noindent\hspace*{1\zw}\begin{tblr}{colspec={*{2}{Q[l,m,wd=62mm]}}, colsep=0pt, rowsep=2pt}
（1）6 × 7 ＝ \bkkakikomi & （2）7 × 8 ＝ \bkkakikomi \\
（3）8 × 9 ＝ \bkkakikomi & （4）9 × 6 ＝ \bkkakikomi \\
（5）7 × 7 ＝ \bkkakikomi & （6）8 × 3 ＝ \bkkakikomi \\
\end{tblr}
\bkpans{（1）42\quad（2）56\quad（3）72\quad（4）54\quad（5）49\quad（6）24}

\bkproblem いちごが 9こずつ 4さらで 36こ。 5さらなら なんこ？
\bkshikiran
\bkpans{\mbox{9 × 5 ＝ 45}\quad 45こ}

\bkproblem □に あう かずを かこう。\par\vspace{2pt}
\noindent\hspace*{1\zw}（1）\mbox{8 × 3 ＝ 3 ×} \bkkakikomi\hspace{3\zw}（2）\mbox{7 × 6 ＝} \bkkakikomi\enspace \mbox{× 7}
\bkpans{（1）8\quad（2）6}

\bkproblem □に \R{入}{はい}る かずを、 \R{表}{ひょう}で さがそう。\par\vspace{2pt}
\noindent\hspace*{1\zw}（1）\mbox{7 ×} \bkkakikomi\enspace ＝ 56\hspace{3\zw}（2）\mbox{9 ×} \bkkakikomi\enspace ＝ 27
\bkpans{（1）8\quad（2）3}

\bkproblem \R{九九}{く|く}の \R{表}{ひょう}の いちぶです。 あいて いる マスに かずを かこう。
\begin{bkfigure}
\begin{tikzpicture}[x=1mm, y=1mm]
  \def\s{9}
  \fill[bkTheme!12] (\s,0) rectangle ({4*\s},-\s) (0,-\s) rectangle (\s,{-4*\s});
  \foreach \c/\t in {1/6,2/7,3/8} {\node[font=\sffamily\bfseries\small, text=bkOrange!85!black] at ({(\c+0.5)*\s},{-\s/2}) {\t};}
  \foreach \r/\t in {1/7,2/8,3/9} {\node[font=\sffamily\bfseries\small, text=bkRoute2!85!black] at ({\s/2},{-(\r+0.5)*\s}) {\t};}
  \foreach \r/\c/\t in {1/1/42,1/2/49,2/2/56,3/1/54,3/3/72} {\node[font=\small] at ({(\c+0.5)*\s},{-(\r+0.5)*\s}) {\t};}
  \foreach \r/\c in {1/3,2/1,2/3,3/2} {\draw[bkTheme, line width=0.9pt, rounded corners=1mm, fill=white] ({(\c+0.5)*\s-3.4},{-(\r+0.5)*\s-2.7}) rectangle ({(\c+0.5)*\s+3.4},{-(\r+0.5)*\s+2.7});}
  \draw[bkGray!45, line width=0.45pt] (0,{-4*\s}) grid[step=\s] ({4*\s},0);
  \draw[bkInk!70, line width=0.9pt] (0,{-4*\s}) rectangle ({4*\s},0) (\s,0) -- (\s,{-4*\s}) (0,-\s) -- ({4*\s},-\s);
  \node[font=\sffamily\bfseries\small, text=bkTheme] at ({\s/2},{-\s/2}) {×};
\end{tikzpicture}
\end{bkfigure}
\bkpans{\mbox{7 × 8 ＝ 56}、 \mbox{8 × 6 ＝ 48}、 \mbox{8 × 8 ＝ 64}、 \mbox{9 × 7 ＝ 63}}

\clearpage
\bkproblemgroup{チャレンジ}{1}

\bkproblem \mbox{9 × 7} を、 5のだんと 4のだんで もとめよう。\par\vspace{2pt}
\noindent\hspace*{1\zw}\mbox{5 × 7 ＝} \bkkakikomi\enspace\mbox{4 × 7 ＝} \bkkakikomi\enspace\mbox{9 × 7 ＝} \bkkakikomi
\bkpans{\mbox{5 × 7 ＝ 35}、 \mbox{4 × 7 ＝ 28}、 \mbox{35 ＋ 28 ＝ 63}\quad \mbox{9 × 7 ＝ 63}}

\bkproblem \R{九九}{く|く}の \R{表}{ひょう}で、 こたえが 24に なる マスを ぜんぶ さがそう。\par\vspace{2pt}
\noindent\hspace*{1\zw}\bkkakikomi[20mm] \hspace{1\zw}\bkkakikomi[20mm] \hspace{1\zw}\bkkakikomi[20mm] \hspace{1\zw}\bkkakikomi[20mm]
\bkpans{\mbox{3 × 8}、 \mbox{4 × 6}、 \mbox{6 × 4}、 \mbox{8 × 3} の 4つ}

\bkproblem \mbox{8 × 10} は いくつ？ \R{表}{ひょう}には ないけれど、 ふえかたの きまりで もとめよう。
\bkshikiran[3\zw]
\bkpans{\mbox{8 × 9 ＝ 72}、 \mbox{72 ＋ 8 ＝ 80}\quad 80}

\begin{furikaeri}\bkqlrag
\begin{checklist}
\item \R{九九}{く|く}の \R{表}{ひょう}で、 こたえを \R{見}{み}つけられる。\hfill{\footnotesize\sffamily\color{bkTheme}→\pageref{bk:1-13:yomikata}ページ}
\item \R{表}{ひょう}の きまりを \R{見}{み}つけられる。\hfill{\footnotesize\sffamily\color{bkTheme}→\pageref{bk:1-13:fuekata}ページ}
\item \R{入}{い}れかえても、 こたえは \R{同}{おな}じと わかる。\hfill{\footnotesize\sffamily\color{bkTheme}→\pageref{bk:1-13:irekae}ページ}
\end{checklist}
チェックが つかなかったら、 →の ページに もどろう。
\end{furikaeri}

\bkfukidashi{sora}{\R{表}{ひょう}を よこや ななめから \R{見}{み}て、 ほかにも きまりを さがして みよう！}
\clearpage
